% ===================================================================
%  TÁBLA Uimh. 7 — An Sráidbhaile sa Gheimhreadh. — An Fheirm san Earrach.
%  Traduction 2026 d'apres texto/io, controlee sur texto/fr, texto/en
%  et texto/es. Consigne : voir texto/ga/10-tabla-01.tex.
%
%  OUVERTURE DE LA DEUXIEME SERIE : « AN DARA SRAITH ». Noter que
%  « sraith » ne mute pas apres « an dara », alors qu'il mutait apres
%  « an chéad » : c'est pourquoi SERIO doit ecrire « sh?raith ».
%
%  LE TABLEAU DE LA FERME EST CELUI OU L'IRLANDAIS EST LE PLUS
%  LUI-MEME : « bó » la vache, « caora » la brebis, « gabhar » la
%  chevre, « capall » le cheval, « muc » le cochon, « cearc » la
%  poule, « tobar » le puits, « scioból » la grange, « aoire » le
%  berger. Pas un emprunt. C'est le meme resultat qu'en basque et en
%  ukrainien, et la meme raison.
% ===================================================================

%%K t07-noto-1 noto
\VUnotes{143.2mm}{%
(*) Féach sonraí an dinnéir i dtáblaí 6 agus 13.%
}

%%K t07-apar-1 apar
\VUsaut{4.0mm}
\VUcentre{13.2pt}{90}{AN DARA SRAITH}
\VUsaut{3.0mm}
\VUfilet{16mm}
\VUsaut{5.6mm}
\VUtitre{15.0pt}{60}{TÁBLA Uimh. 7}
\VUsaut{3.0mm}

%%K t07-tit-1 sub
\VUcentre{12.0pt}{0}{\textit{An Chéad Radharc.}}
\VUsaut{3.0mm}

%%K t07-tit-2 sub
\VUpk{10.2pt}{40}{An Sráidbhaile sa Gheimhreadh.}
\VUsaut{4.0mm}

%%K t07-c1-01-1 p
1. --- Ar shaoire na hAthbhliana chuaigh mé le m'athair go Baile Nua, sráidbhaile
ina raibh gnó aige.

%%K t07-c1-02-1 p
2. --- Chuaigh muid sa \VUgras{chóiste} \textsuperscript{(11)}, a théann idir an
baile mór agus an sráidbhaile seo; ach ó bhí sé an-fhuar, níorbh aoibhinn ar chor
ar bith an turas i gcarr chomh míchompordach.

%%K t07-c1-03-1 p
3. --- Bhíomar an-sásta, mar sin, nuair a stop an \VUgras{tiománaí} a chóiste sa
chearnóg, os comhair \VUgras{tábhairne} \textit{«An Béar Bán»}
\textsuperscript{(2)}. Bhí \VUgras{fear an tábhairne} \textsuperscript{(4)} ag
fanacht linn ar an tairseach, ag caitheamh a \VUgras{phíopa} go socair.

%%K t07-c1-03-2 p
Thug sé cuireadh dúinn teacht isteach, agus níorbh éigean mé a iarraidh faoi dhó
chun suí os comhair na tine fíniúna a bhí ag cnagadh sa teallach. Ansin ba chuma
liom faoin \VUgras{ngaoth aduaidh} ghéar, a chuir an \VUgras{fhógra}
\textsuperscript{(3)} sean ag díoscán agus a d'iompair i ndlaoithe dubha an
\VUgras{deatach} \textsuperscript{(1)} a bhí ag teacht as an simléar.

%%K t07-c1-04-1 p
4. --- Ba é am an dinnéir é. Gan mhoill thugamar an urraim chuí don bhia simplí ach
blasta a tugadh dúinn (*).

%%K t07-tit-3 sub
\VUpk{10.2pt}{40}{Roinnt cuairteanna.}
\VUsaut{3.0mm}

%%K t07-c1-05-1 p
5. --- Tar éis an dinnéir chuaigh mé le m'athair, atá ina ghníomhaire árachais,
chuig a chuid cliant. Ar dtús thugamar cuairt ar an \VUgras{ngabha}
\textsuperscript{(5)}, atá ina chónaí in aice leis an tábhairne. Bhí sé ag cur crúite
ar chapall. Chuir neart a chúntóra iontas orm, a choinnigh \VUgras{crúb} an chapaill
go daingean lena naprún leathair, agus an gabha ag bualadh na \VUgras{crúite}
\textsuperscript{(6)} le buillí troma casúir.

%%K t07-c1-06-1 p
6. --- Ansin chuaigh muid chuig \VUgras{gréasaí} an tsráidbhaile
\textsuperscript{(7)}, chuig Séamas críonna. Cé gur \VUgras{buatais} iontach a bhí
mar fhógra aige, ní raibh sé ach ag deisiú péire sean de bhróga caite.

%%K t07-c1-06-2 p
Dúirt an seanfhear linn agus é ag gáire: «Sa bhaile mór ba „ghréasaí” mé, agus ní
raibh custaiméirí agam. Anseo is „deisitheoir” mé, agus tá obair go leor agam.»

%%K t07-c1-07-1 p
7. --- Labhair sé linn faoina chomharsa an \VUgras{bearbóir}
\textsuperscript{(8)}, a bhfuil a chuid custaiméirí go léir míshásta. Bíonn a
\VUgras{rásúir} maolaithe i gcónaí, agus ní ghearrann a \VUgras{shiosúr} riamh mar
is ceart. Ach ó tharla gurb é an t-aon bhearbóir sa tsráidbhaile é, níl aon rogha
eile ag na daoine ach a mbearradh agus a lomadh a fháil uaidh. Ina theannta sin is
duine sprionlaithe agus mífhial é.

%%K t07-c1-08-1 p
8. --- Ar an dea-uair, níl na \VUgras{comharsana} eile cosúil leis. Is duine maith,
dícheallach agus cabhrach é an \VUgras{roithleoir} \textsuperscript{(9)}, mar
shampla. Is pléisiúr an charr a thabhairt dó lena dheisiú. Bíonn a chuid oibre
déanta go glan i gcónaí.

%%K t07-c1-09-1 p
9. --- Níor ghearán Séamas críonna faoin \VUgras{diallaitiúir}
\textsuperscript{(10)} ach oiread, a gcabhraíonn sé leis uaireanta na
\VUgras{húmacha} a fhuáil, nuair nach bhfulaingíonn an obair moill.

%%K t07-c1-09-2 p
D'fhiafraigh m'athair de faoin teaghlach. «Tá gach duine slán. Tá m'iníon mhic ar
\VUgras{scoil} \textsuperscript{(12)}. Tá an \VUgras{mhúinteoir}
\textsuperscript{(13)} an-sásta léi; is \VUgras{dalta} maith í
\textsuperscript{(14)}. Ní hionann agus an dailtín seo, mo mhac; níl ina cheann ach
cluichí. Ní féidir leis an \VUgras{múinteoir} \textsuperscript{(15)} rud ar bith a
chur á fhoghlaim aige.»

%%K t07-tit-4 sub
\VUsaut{3.0mm}
\VUpk{10.2pt}{40}{Cúlchaint an tsráidbhaile.}
\VUsaut{3.0mm}

%%K t07-c1-10-1 p
10. --- Ansin d'inis an seanfhear dúinn nuacht an tsráidbhaile. Bhí siad ar tí scoil
nua a thógáil, mar tá an ceann atá sa \VUgras{halla baile} imithe róchúng. Níorbh
obair mhór é, mar ba leor cuid de ghairdín an \VUgras{mhuilleora} a cheannach.
Bheadh sé an-áisiúil don fhear bocht. Ó thosaigh na sioctha níorbh fhéidir le
\VUgras{clocha} an \VUgras{mhuilinn} \textsuperscript{(16)} an t-arbhar a
mheilt, mar tá an \VUgras{roth} \textsuperscript{(17)} greamaithe ag an
\VUgras{oighear} \textsuperscript{(25)} ar an \VUgras{sruthán}
\textsuperscript{(23)}.

%%K t07-c1-11-1 p
11. --- «Ó tharla go bhfuilimid ag caint ar thógálacha: an bhfaca sibh ár
\VUgras{séipéal} nua \textsuperscript{(18)}? Tá sé an-phictiúrtha faoina chaipín
sneachta. Suíonn na \VUgras{préacháin} \textsuperscript{(19)}, nach n-aithníonn a
sean-chloigtheach a thuilleadh, go gruama ar an gcrann mór sa
\VUgras{reilig} \textsuperscript{(20)}.»

%%K t07-c1-12-1 p
12. --- Chuir trácht na reilige i gcuimhne don ghréasaí na daoine a bhí ar aithne ag
m'athair agus a fuair bás i rith na bliana. «Nach iomaí \VUgras{uaigh}
\textsuperscript{(21)}, a dhuine uasail, atá curtha leis na sean-chinn faoin
\VUgras{gcufróg} \textsuperscript{(22)}! Thóg an bás ár sean-nótaire. Bhí an
sráidbhaile ar fad ag a \VUgras{shochraid}.»

%%K t07-c1-13-1 p
13. --- «Sin é a chomharba ag scátáil ar an sruthán. Tháinig sé an lá faoi dheireadh
chun go ndeisínn strapa a \VUgras{scáta} \textsuperscript{(26)}, a bhí briste.»

%%K t07-c1-14-1 p
14. --- «Agus thall ansin, ar bhruach an tsrutháin, an \VUgras{maor páirce} nua
\textsuperscript{(27)}. Is iar-shaighdiúir é agus is scanradh dhailtíní an
tsráidbhaile é. Scaipeann siad chomh luath agus a fheiceann siad a
\VUgras{loirgníní} \textsuperscript{(29)} agus a \VUgras{chaipín}
\textsuperscript{(28)}.»

%%K t07-c1-15-1 p
15. --- «Á! sin é Iósaef críonna ag tabhairt \VUgras{cairte brosna}
\textsuperscript{(30)} don bháicéir. --- Go deimhin, arsa m'athair, aithním a
fhoireann \VUgras{miúileanna} \textsuperscript{(31)}. Caithfidh mé labhairt leis, agus
bainfidh mé leas as seo. Slán agat, a Shéamais chríonna.»

%%K t07-c1-16-1 p
16. --- Díreach agus muid ag imeacht, bhí páistí an tsráidbhaile ag teacht amach ón
scoil agus ag caitheamh iad féin ar an \VUgras{sleamhnán oighir}
\textsuperscript{(33)}, agus baol ann go dtitfidís agus go mbrisfidís rud éigin sa
\VUgras{titim} \textsuperscript{(34)}. Thóg duine acu ón roithleoir a
\VUgras{shleamhnán} \textsuperscript{(32)}, chuir sé comrádaí ann agus rith sé leis.

%%K t07-c1-17-1 p
17. --- Chaith cuid eile iad féin ar an \VUgras{mbáisín sneachta}
\textsuperscript{(37)} in aice leis an \VUgras{droichead} \textsuperscript{(40)}
agus thosaigh siad ag caitheamh liathróidí leis an \VUgras{bhfear sneachta}
\textsuperscript{(39)}, a rinne siad an lá roimhe sin agus ar thug siad hata, píopa
agus mála scoile ar a ghualainn dó.

%%K t07-c1-18-1 p
18. --- Ba chuma leo faoin ngaoth oighreata agus faoin bhfuacht. Ní raibh fiú
\VUgras{scaif} \textsuperscript{(35)} ag a bhformhór, agus chun iad féin a théamh
tar éis an chluiche \VUgras{liathróidí sneachta} \textsuperscript{(38)}, ba leor
dorn de \VUgras{chastáin} the \textsuperscript{(42)}, ceannaithe ó Jacqueline
chríonna, an \VUgras{bhean díola}, atá ar crith leis an bhfuacht faoina
\VUgras{seál} róthanaí \textsuperscript{(43)}.

%%K t07-tit-5 sub
\VUsaut{3.0mm}
\VUcentre{12.0pt}{0}{\textit{An Dara Radharc.}}
\VUsaut{3.0mm}

%%K t07-tit-6 sub
\VUpk{10.2pt}{40}{An Fheirm san Earrach.}
\VUsaut{4.0mm}

%%K t07-c2-01-1 p
1. --- Le gach áthas an gheimhridh, fáiltímid le lúcháir mhór roimh fhilleadh an
\VUgras{Earraigh}.

%%K t07-c2-01-2 p
Nach taitneamhach an pictiúr é clós na feirme tráthnóna, i mí na Bealtaine!

%%K t07-c2-02-1 p
2. --- Ag bun na spéire tá an \VUgras{ghrian} \textsuperscript{(1)} ag dul faoi go
mall, ag báthadh na \VUgras{bpáirceanna} \textsuperscript{(3)} le \VUgras{gathanna}
corcra \textsuperscript{(2)}. Tá an \VUgras{treabhdóir} dícheallach
\textsuperscript{(6)} ag deifriú chun a \VUgras{ghort} \textsuperscript{(4)} a
threabhadh.

%%K t07-c2-02-2 p
Lena \VUgras{chéachta} \textsuperscript{(5)}, atá á tharraingt ag \VUgras{capall}
láidir \textsuperscript{(7)}, tarraingíonn sé an \VUgras{eitre}
\textsuperscript{(8)}, ina scaipfidh an \VUgras{síoladóir} \textsuperscript{(9)} le
doirne an \VUgras{síol} a thabharfaidh diasa óir.

%%K t07-c2-03-1 p
3. --- Bhain na \VUgras{buanaithe} \textsuperscript{(11)} lena speala féar na
móinéire, thriomaigh na hiompóirí an \VUgras{féar} \textsuperscript{(10)}, á
iompú lena \VUgras{bpíce} \textsuperscript{(12)} agus lena rácaí, agus seo iad ag
filleadh abhaile.

%%K t07-c2-03-2 p
Siúlann cuid acu os comhair an chairr throim, atá á tharraingt ag daimh, agus na
huirlisí ar a nguaillí. Tá cuid eile, atá níos leisciúla, socraithe go
compordach ar an bhféar cumhra.

%%K t07-c2-04-1 p
4. --- Ag dul thar bráid, beannaíonn siad go cairdiúil don \VUgras{gharraíodóir}
\textsuperscript{(16)}, atá ag saothrú sa \VUgras{ghairdín}
\textsuperscript{(13)}, ag bearradh na \VUgras{gcrann torthaí} agus ag beangú na
\VUgras{bpiorraí} \textsuperscript{(14)}. Tá na \VUgras{plumaí}
\textsuperscript{(15)} faoi bhláth, agus sa \VUgras{ghlasraígh}
\textsuperscript{(17)} atá leasaithe go maith tá na glasraí, an cabáiste agus an
leitís go breá cheana.

%%K t07-c2-04-2 p
Ar an bhfál íseal osclaíonn \VUgras{péacóg} álainn \textsuperscript{(18)} a
heireaball, chun a cuid cleití iontacha a thaispeáint.

%%K t07-tit-7 sub
\VUpk{10.2pt}{40}{Clós na feirme.}
\VUsaut{3.0mm}

%%K t07-c2-05-1 p
5. --- Tá doirse an \VUgras{stábla} \textsuperscript{(19)} ar oscailt, agus
feictear an mhainséar agus \VUgras{leaba} \textsuperscript{(23)} na \VUgras{lárach}
\textsuperscript{(21)}, a bhain an \VUgras{giolla capaill}
\textsuperscript{(20)} den chrann díreach anois. Siúlann an \VUgras{searrach}
\textsuperscript{(22)}, chomh greannmhar sin ar a chosa fada, i ndiaidh a mháthar,
ag pramsáil.

%%K t07-c2-06-1 p
6. --- In aice leis an \VUgras{tobar} \textsuperscript{(29)} tagann na
\VUgras{ba} \textsuperscript{(25)} chun ól as \VUgras{umar} adhmaid
\textsuperscript{(26)}, atá ina n-áit uisce. Baineann an \VUgras{lao}
\textsuperscript{(24)} leas as an deis chun diúl, agus déanann an \VUgras{damh}
\textsuperscript{(27)}, agus boladh maith an fhodair aige, géimneach ghlórach agus
ardaíonn sé go bródúil a cheann lena \VUgras{adharca} cumhachtacha
\textsuperscript{(28)}.

%%K t07-c2-07-1 p
7. --- Tá gach duine ag obair. Coinníonn an \VUgras{bhean oibre}
\textsuperscript{(32)} \VUgras{téad} an tobair \textsuperscript{(31)}. Tarraingíonn
sí uisce leis an \VUgras{mbuicéad} \textsuperscript{(30)}, chun deoch a thabhairt do
na hainmhithe. In aice leis an \VUgras{scioból} \textsuperscript{(33)} tá fear ag
cruinniú an tuí atá scaipthe, agus os comhair dhoras na \VUgras{feirme}
\textsuperscript{(34)} tá \VUgras{sníomhaí} chríonna \textsuperscript{(35)} lena
tuirne agus lena \VUgras{coigeal} ag sníomh an \VUgras{lín} nó an
\VUgras{chnáib}, mar a bhíodh sna seanlaethanta.

%%K t07-tit-8 sub
\VUsaut{3.0mm}
\VUpk{10.2pt}{40}{An Feirmeoir.}
\VUsaut{3.0mm}

%%K t07-c2-08-1 p
8. --- Stiúrann an \VUgras{feirmeoir} \textsuperscript{(36)} an obair seo go léir
agus tugann sé treoracha dá dhaoine. Ordaíonn sé an \VUgras{bhráca}
\textsuperscript{(40)} agus an \VUgras{piocóid} \textsuperscript{(39)} a chur faoin
sciathán, atá fágtha in aice le \VUgras{both} \textsuperscript{(38)} an
\VUgras{mhadra} \textsuperscript{(37)}.

%%K t07-c2-09-1 p
9. --- Scanraíonn a chuid glaonna an \VUgras{colúr} \textsuperscript{(41)}, a
eitlíonn ar a dhícheall chuig an \VUgras{gcolúrlann} \textsuperscript{(42)}, agus na
\VUgras{cearca} \textsuperscript{(43)}, a dheifríonn ar ais chuig an
\VUgras{gcró} \textsuperscript{(44)}.

%%K t07-c2-10-1 p
10. --- Os comhair na \VUgras{caoraigheann} \textsuperscript{(48)} tá an
\VUgras{t-aoire} críonna \textsuperscript{(45)} ag cruinniú a \VUgras{thréada}
\textsuperscript{(49)}. Agus a \VUgras{bhachall} \textsuperscript{(46)} ina láimh,
nach scarann leis, mar nach scarann a \VUgras{bhlús} tréigthe
\textsuperscript{(47)} ach oiread, tiomáineann sé isteach sa loca na
\VUgras{reithí} \textsuperscript{(50)}, na \VUgras{caoirigh}
\textsuperscript{(53)}, na \VUgras{huain} \textsuperscript{(54)}, na
\VUgras{gabhair} \textsuperscript{(51)} agus na \VUgras{meannáin}
\textsuperscript{(52)}. Siúlann a \VUgras{mhadra} dílis \textsuperscript{(55)},
Argus, ar deireadh agus deifríonn sé lena thafann na cinn atá ar gcúl.

%%K t07-c2-11-1 p
11. --- Ní chuireann an torann agus an fuadar seo go léir isteach ar shuaimhneas na
\VUgras{bó bhainne} maithe \textsuperscript{(56)}, a bhíonn ag baint gliogar as a
\VUgras{clog} \textsuperscript{(57)} agus í á crú os comhair dhoras na
\VUgras{bóithe} \textsuperscript{(58)}.

%%K t07-c2-12-1 p
12. --- Dealraíonn sí a thuiscint go gcaithfidh sí bainne a thabhairt, chun go
bhféadfaidh \VUgras{bean an fheirmeora} \textsuperscht{(60)} an \VUgras{t-im} a
mhaistreadh sa \VUgras{chuinneog} \textsuperscript{(61)} nó \VUgras{cáiseanna}
den scoth a dhéanamh \textsuperscript{(64)}, a thriomaíonn ar chlár curtha trasna an
\VUgras{dabhaigh} \textsuperscript{(63)}.

%%K t07-c2-13-1 p
13. --- Coinnítear an \VUgras{déirí} \textsuperscript{(59)} ar fad glan ag an
\VUgras{mbuachaill aimsire} \textsuperscript{(65)}, a nigh na \VUgras{cannaí}
\textsuperscript{(62)} díreach anois sa bhuicéad mór atá á iompar aige.

%%K t07-tit-9 sub
\VUsaut{3.0mm}
\VUpk{10.2pt}{40}{Clós na n-éan.}
\VUsaut{3.0mm}

%%K t07-c2-14-1 p
14. --- Ina bhróga adhmaid troma siúlann sé go místuama go leor agus scanraíonn sé
dá bhrí sin an \VUgras{turcaí} \textsuperscript{(68)}, na \VUgras{lachain}
\textsuperscript{(66)}, na \VUgras{bardail} \textsuperscript{(67)} agus na
\VUgras{géanna} \textsuperscript{(69)}, a osclaíonn a \VUgras{ngoba}
\textsuperscript{(70)} go leathan agus a thógann gleo imníoch.

%%K t07-c2-14-2 p
Ardaíonn an \VUgras{coileach} \textsuperscript{(71)}, atá níos cathaithí, a
\VUgras{chíor} dhearg \textsuperscript{(72)}, croitheann sé cleití a
\VUgras{eireabaill} \textsuperscript{(73)} agus scaoileann sé «cuc-a-dú-dal-dú»
glórach.

%%K t07-c2-15-1 p
15. --- Bíonn an \VUgras{chearc ghoir} \textsuperscript{(74)} ag scríobadh sa
\VUgras{charn aoiligh} \textsuperscript{(81)} lena \VUgras{sicíní}
\textsuperscript{(75)}. Ritheann an \VUgras{banbh} \textsuperscript{(78)} chuig a
mháthair, chuig an \VUgras{gcráin} ollmhór \textsuperscript{(77)}, atá le feiceáil
in aice leis an gcró.

%%K t07-c2-16-1 p
16. --- Ina gcliabháin itheann na \VUgras{coiníní} \textsuperscript{(79)} go santach
an \VUgras{féar} bog \textsuperscript{(80)}, a thugann iníon an fheirmeora chucu. Is
mór ag Eibhlín a cuid coiníní, agus gach lá, tar éis di na \VUgras{huibheacha}
úra \textsuperscript{(76)} a thógáil sa chró, tagann sí ar cuairt chuig a cairde
beaga.
\VUsaut{6.0mm}
\VUfilet{22mm}
